\section{Two-dimensional Kramers-Wannier duality and fusion 2-categorical symmetry}
Having established the tensor network framework of categorical symmetries and dualities in one-dimensional lattice models, we now proceed to the two-dimensional Kramers-Wannier duality and the associated higher symmetry structures described by \emph{fusion 2-categories}~\cite{Wegner1971,Haegeman2014a,Douglas2018,Delcamp2023,Decoppet2024}. Since fusion 2-categories are considerably more intricate structures than their 1-categorical counterparts, we keep the following discussion as concrete as possible. We relegate a sketch of the general framework of fusion 2-categorical symmetries and their dualities to the following section. We begin by constructing the intertwining PEPO that implements the two-dimensional Kramers-Wannier duality.
\subsection{Two-dimensional Kramers-Wannier duality}
Let us consider for concreteness the two-dimensional \emph{triangular lattice}, endowed with qubits on the vertices $\msf V$. Similar to the one-dimensional case, we will study the fate of the two-dimensional transverse field Ising model under Kramers-Wannier duality. In two dimensions, the Hamiltonian reads
\begin{equation}\label{eq:2dTFIM}
	H = -J\sum_\msf e Z_{s(\msf e)}Z_{t(\msf e)} - \lambda \sum_\msf v X_\msf v.
\end{equation}
Herein, $s$ and $t$ are functions that assign to every edge their \emph{source} and \emph{target}, vertices. Let us assume \emph{periodic} boundary conditions, so that effectively the underlying topology is that of the two-torus. Evidently, the Hamiltonian commutes with a global $\mbb Z_2$ spin-flip symmetry generated by $\prod_\msf v X_\msf v$. In the context of higher-form symmetries, such a symmetry operator is called a \emph{0-form} symmetry operator or \emph{surface operator} since the operator extends over all of space, or thus a submanifold of  \emph{codimension} 0 with respect to the spatial dimensionality. Notably, one could add to \cref{eq:2dTFIM} additional terms that commute with the 0-form $\mbb Z_2$ symmetry without altering much of the following discussion~\cite{Delcamp2023}.

\bigskip\noindent In analogy with the one-dimensional Kramers-Wannier duality operator \cref{eq:KWMPO}, its two-dimensional counterpart admits a PEPO representation of bond dimension two that can be depicted as~\cite{Haegeman2014a}
\begin{equation}\label{eq:KWPEPO}
	\inputtikz{KWPEPO}\qquad\,,
\end{equation}
where the constituent tensors were defined in \cref{eq:TCTensors}. By pulling the Hamiltonian terms through this intertwiner and thereby making use of the symmetry properties listed in \cref{eq:TCTensorsSyms}, one finds the following mapping:
\begin{equation}\label{eq:2dKWMapping}
	Z_{s(\msf e)}Z_{t(\msf e)} \mapsto Z_\msf e, \qquad X_\msf v \mapsto \prod_{\msf e|\msf v} X_\msf e,
\end{equation}
where $\msf e|\msf v$ stands for the collection of all edges $\msf e\in\msf E$ adjacent to $\msf v$. Collecting everything, the dual Hamiltonian is thus found to be
\begin{equation}\label{eq:2dDualTFIM}
	H^\vee = -J\sum_\msf e Z_\msf e - \lambda \sum_\msf v \prod_{\msf e|\msf v} X_\msf e.
\end{equation}
This is however not the complete story, since it can be easily verified that the dual model obeys the kinematical constraints 
\begin{equation}\label{eq:TCKineConst}
	\prod_{\msf e\in\ell} Z_\msf e \overset{!}{=} \mbb 1
\end{equation}
for every closed loop $\ell$ on the triangular lattice, including the non-contractible ones. Similarly as in one dimension, this signals in particular that the dual model is in the even sector of a dual \emph{1-form} symmetry that we discuss momentarily. In conclusion, the dual model can thus be thought of as a toric code in a transverse field, wherein the plaquette terms are enforced kinematically rather than dynamically, and where the terms in the second sum of \cref{eq:2dDualTFIM} are the vertex terms.
\subsection{Dual symmetry structure}
Let us proceed by unpacking the full symmetry structure of the dual model $H^\vee$. First, for every closed loop $\ell$, operators of the form
\begin{equation}
	\mbb U[\ell] := \prod_{\msf e\in\ell} Z_\msf e
\end{equation}
are immediately verified to commute with $H^\vee$. Moreover, they are \emph{topological} by virtue of the kinematical constraints \cref{eq:TCKineConst}, and in the context of gauge theories referred to as \emph{Wilson loops}~\cite{Wegner1971,Wilson1974,Kogut1979,Kogut1982}. Taken together, this means that the collection of loop operators $\mbb U[\ell]$ for all closed loops $\ell$ generate a \emph{1-form} $\Rep(\mbb Z_2)$ symmetry. \footnote{Let us note at this point that in general, the gauging of a $(d+1)$-dimensional theory with a $q$-form symmetry is expected to result in a dual model with a $(d-q)$-form symmetry~\cite{Gaiotto2014}, see also \cref{app:codes}.}

Recently, it has been appreciated that besides the Wilson loops, the dual model hosts \emph{surface operators} that are obtained by \emph{condensing} the loop operators on the spatial manifold~\cite{Kong2014,Gaiotto2019,JohnsonFreyd2020}. Surface operators that admit this interpretation are commonly referred to in the literature as \emph{condensation defects} and are often thought of as resulting from a \emph{higher gauging} procedure~\cite{Choi2022,Choi2024,Delcamp2021,Lin2022,Roumpedakis2022}. More concretely, by this it is meant that the surface operator results from summing over a \emph{sufficiently fine mesh} of Wilson lines~\cite{Bhardwaj2017}. By making use of the topological invariance and fusion rules of the loop operators, this network can be reduced to the following \emph{minimal} network of Wilson loops:
\begin{equation}\label{eq:CD}
	\mbb U^{\{\mbb 1\}} := \frac{1}{2} \bigg( \sum_{s_x\in\{0,1\}}\mbb U[\ell_x]^{s_x} \bigg) \cdot \bigg( \sum_{s_y\in\{0,1\}} \mbb U[\ell_y]^{s_y} \bigg),
\end{equation}
where $\ell_x$ and $\ell_y$ denote representative loops along the two non-contractible cycles of the two-torus. The notation $\mbb U^{\{1\}}$ for this condensation defect is motivated below.

From its definition \cref{eq:CD} we immediately infer the following key properties. First, in stark contrast with the 0-form $\mbb Z_2$ symmetry of the original model, $\mbb U^{\{\mbb 1\}}$ acts trivially on any local operator. This follows directly from the topological invariance of the 1-form symmetry operators appearing in the definition \cref{eq:CD}. An easy computation also shows that $\mbb U^{\{\mbb 1\}}$ multiplies according to $\mbb U^{\{\mbb 1\}}\cdot\mbb U^{\{\mbb 1\}}=2 \cdot \mbb U^{\{\mbb 1\}}$, so that in particular it does not admit a multiplicative inverse. \footnote{In fact, a more careful computation reveals that on a generic closed and oriented spatial manifold $\Sigma$ the fusion rule of $\mbb U^{\{\mbb 1\}}$ reads $\mbb U^{\{\mbb 1\}}\cdot \mbb U^{\{\mbb 1\}} = Z_{\mbb Z_2}(\Sigma) \cdot \mbb U^{\{\mbb 1\}}$, with $Z_{\mbb Z_2}(\Sigma)=2^{b_1-b_0}$ the partition function of pure $\mbb Z_2$ gauge theory on $\Sigma$ and $b_{0,1}$ the respective \emph{Betti numbers} of $\Sigma$. Hence, the fusion rules of this symmetry operator are \emph{topology dependent}~\cite{Roumpedakis2022,Delcamp2023}.} Finally, note that all Wilson loops condense on this surface in the sense that $\mbb U[\ell]\cdot\mbb U^{\{\mbb 1\}} = \mbb U^{\{\mbb 1\}}= \mbb U^{\{\mbb 1\}}\cdot\mbb U[\ell]$.

\bigskip\noindent Using the tensor network calculus developed in refs.~\cite{Delcamp2021,Delcamp2023} and \cref{app:TwistedGauging}, a PEPO representation of $\mbb U^{\{\mbb 1\}}$ can be constructed. Defining the tensor
\begin{equation}
	\inputtikz{CD_PEPO_tensor_1} \,\, := \,\, \inputtikz{CD_PEPO_tensor_2}\,\, ,
\end{equation}
this PEPO is depicted as:
\begin{equation}
	\inputtikz{CD}\qquad\,.
\end{equation}
Remark the striking similarity with the $\mc R=\Vect_{\mbb Z_2}$ tensor network representation  of the toric code ground state constructed in \cref{sec:TC}. Indeed, in a very loose but not meaningless sense, one can think of the condensation operator as the toric code ground state in which the physical degrees of freedom are `doubled' so as to turn it into a non-invertible surface operator.

The similarity with the toric code ground state is more than an analogy in the sense that the condensation defect can be decorated with topological lines labelled by elements in $\mbb Z_2$. In the language of gauge theories, such lines are referred to as \emph{'t Hooft lines}. Very concretely, one constructs for any loop $\ell^\vee$ on the \emph{dual} lattice, the non-trivial 't Hooft line by inserting a string of Pauli-$X$ operators on the virtual level along $\ell^\vee$, cf. \cref{eq:TC_1_string}. For example:
\begin{equation}\label{eq:CD_string}
		\inputtikz{CD_string}\qquad\,.
\end{equation}
Topological invariance follows directly from the symmetry properties of the local tensors, as usual.

\bigskip\noindent Let us recapitulate what we have found so far. We have argued that the dual model $H^\vee$ enjoys 1-form $\Rep(\mbb Z_2)$ symmetry. By condensing these Wilson loops, one obtains a condensation defect $\mbb U^{\{\mbb 1\}}$ that itself hosts 't Hooft lines organised in $\Vect_{\mbb Z_2}$. By similar methods, one can construct a topological interface between $\mbb U^{\{\mbb 1\}}$ and the trivial identity surface operator. Formally, these structures combine into a \emph{fusion 2-category} denoted $\TRep(\mbb Z_2)$. Notably, $\TRep(\mbb Z_2)$ can be thought of as a \emph{categorification} of the category $\Rep(\mbb Z_2)$, which is the symmetry structure one expects after performing the Kramers-Wannier duality on a $\Vect_{\mbb Z_2}$-symmetric model. And indeed, as seen below, the 0-form $\mbb Z_2$ symmetry of the Ising model provides a concrete realisation of the symmetry 2-category $\TVect_{\mbb Z_2}$. Remarkably, unlike in the one-dimensional setting, the original and dual symmetry are not monoidally equivalent. This prevents us from constructing a model \emph{self-dual} under the Kramers-Wannier duality. In turn, this raises the question whether other duality transformations exist that give rise to (2+1)-dimensional self-dual models. The answer is affirmative. In \cref{app:TwistedGauging}, we consider a more refined duality transformation that amounts to the composition of a \emph{gauging} and \emph{twisted gauging} procedure that admits self-dual models.

Before unpacking at a more abstract level fusion 2-categorical symmetry structures, let us discuss the \emph{topological sectors} of the models from this section and study their permutation under the duality transformation.
\subsection{Topological sectors}
In complete analogy with the one-dimensional case, the two-dimensional Ising model admits periodic or antiperiodic boundary conditions along either non-contractible cycle of the torus. Furthermore, the Hilbert space for each choice of boundary conditions decomposes into two $\mbb Z_2$ charge sectors, so that the total number of topological sectors is $2^3=8$.

To identify these sectors in the dual model, we note first that the Wilson loop operators $\mbb U[\ell_x]$, $\mbb U[\ell_y]$ measure the $\mbb Z_2$ \emph{flux} through the corresponding cycles of the torus. This labels four distinct charge sectors with respect to the dual 1-form symmetry. These can be accessed by choosing the boundary conditions of the Ising model. It can also be verified explicitly that the odd $\mbb Z_2$ charge sector of the \emph{original} model can be accessed by changing the sign of one of the vertex terms in \cref{eq:2dDualTFIM}. Heuristically, this can be seen by taking the product of all vertex terms and tracking that under duality this amounts to
\begin{equation}
	\prod_\msf v X_\msf v = -\prod_\msf v \prod_{\msf e|\msf v} X_\msf e = -\mbb 1.
\end{equation}
Changing the sign of one of these vertex terms is sometimes referred to as choosing a \emph{1-form twisted boundary condition}~\cite{Moradi2022,Moradi2023}. 

In summary, we obtain a very similar mapping as in the one-dimensional case, whereby the boundary conditions and charge sectors are interchanged.
\section{Algebraic structure of fusion 2-categories}
Let us now unpack in some more detail what a fusion 2-categorical symmetry entails and develop its graphical calculus in parallel~\cite{Douglas2018,Delcamp2021}.
\subsection{Algebraic structure of fusion 2-categories}
Loosely speaking, a \emph{fusion 2-category} $\fr C$ describes a collection of \emph{topological surfaces} which constitute the object set $\ob(\fr C)$. In what follows we will use \emph{(topological) surfaces} and \emph{objects} interchangeably. The physically inclined reader might want to think about these surfaces as being implemented in a lattice model via PEPOs. Surfaces which cannot be written as the direct sum of other surfaces are said to be \emph{simple}. We write this set as $\mc I_\fr C$. Fusing two simple surfaces and decomposing the result can be depicted as
\begin{equation}
	\inputtikz{2cat_fusion_1}\,\,\,\,  = \bigoplus_{X_3\in\mc I_{\fr C}} N_{X_1X_2}^{X_3} \,\, \inputtikz{2cat_fusion_2}\,\, ,
\end{equation}
for an integer rank-3 tensor $N$. The fusion rules are associative. Among the simple objects, there exists a unique \emph{monoidal unit} denoted by $\mbb 1_\fr C$, or simply $\mbb 1$. Sometimes we will refer to $\mbb 1_\fr C$ as the \emph{trivial} or \emph{transparent} surface. \emph{Topological defect line}s on the interface between generic surfaces $X$ and $X'$ are organised in a \emph{Hom (1-)category} denoted by $\fr C(X,X')$, and we will use the shorthand $\fr C^{X,X'}\equiv\fr C(X,X')$. On the lattice, these might be represented via MPOs, living on the \emph{virtual level} of the PEPOs, akin to those defined in \cref{eq:CD_string}. Topological lines are said to encode the \emph{1-morphisms} of the 2-category. Graphically:
\begin{equation}
	\inputtikz{2cat_1morph_1}\,\, = \inputtikz{2cat_1morph_2}\,\,, 
\end{equation}
for $f\in\ob(\fr C(X,X'))$. Lines between compatible surface operators can be fused `horizontally':
\begin{equation}
	\inputtikz{2cat_1morph_3} \,\, = \,\, \inputtikz{2cat_1morph_4}\,\, .
\end{equation}
It can be shown that the Hom category $\fr C^{X,X}$ for any simple object $X$ naturally forms a \emph{fusion 1-category}. Specialising to the identity object $X=\mbb 1$, it is found that $\fr C^{\mbb 1,\mbb 1}$ is in fact a \emph{braided} fusion category, but the braiding is not usually modular. And remarkably, $\fr C^{X,X'}$ is an indecomposable $(\fr C^{X,X},\fr C^{X',X'})$-\emph{bimodule category} for any pair of simple objects $X,X'$.

Finally, the \emph{morphisms} in $\fr C^{X,X'}$ encode topological \emph{point operators} on the defect lines, and are said to constitute the \emph{2-morphisms} of $\fr C$. These can be composed `vertically', which is pictorially represented as:
\begin{equation}
	\inputtikz{2cat_2morph_1} \,\, = \,\, \inputtikz{2cat_2morph_2}\,\, ,
\end{equation}
where $f,g,h\in\ob(\fr C(X,X'))$, $\varphi\in\fr C^{X,X'}(f,g)$, $\psi\in\fr C^{X,X'}(g,h)$, and $\chi=\psi\circ\varphi$.

Associativity of the tensor product is encoded by natural \emph{associator} 1-isomorphisms $\alpha^{X_1X_2X_3}:(X_1\otimes X_2)\otimes X_3 \overset{\sim}{\rightarrow} X_1\otimes (X_2\otimes X_3)$. The associator satisfies the pentagon equation up to natural 2-isomorphisms called the \emph{pentagonator}~\cite{Douglas2018,Delcamp2021}. 

\bigskip\noindent In stark contrast with the 1-categorical setting, fusion 2-categories thus allow for the existence of non-trivial topological defects on the interface between distinct simple objects, as well as \emph{self-interfaces}. This naturally leads to the definition of a \emph{connected component}, which is an equivalence class of simple objects with the equivalence relation $X\sim X'\iff \fr C(X,X')\neq 0$. Physically, all surfaces that belong to the same connected component are related via \emph{condensation} of an appropriate set of topological lines, as we will exemplify below~\cite{Kong2014,Gaiotto2019,JohnsonFreyd2020,Roumpedakis2022}. Moreover, the surfaces belonging to the identity component are exactly those that can be defined on \emph{open} patches. As such, any fusion 2-category admits a finite direct sum decomposition:
\begin{equation}
	\fr C \simeq \bigoplus_{i=1}^r \fr C_i.
\end{equation}
\subsection{Examples}
\subsubsection{Finite group symmetry: $\TVect_G^\pi$}
Arguably the simplest fusion 2-category is $\TVect_G$, for any finite group $G$. The simple objects coincide with $\Vect_g$, often identified with the group elements, since the fusion rules simply follow the multiplication of $G$: $\Vect_{g_1}\otimes\Vect_{g_2} = \Vect_{g_3}$. For the particular case of $G=\mbb Z_2$, we obtain the symmetry structure of the two-dimensional Ising model introduced and discussed in the previous section. Notably, there are no non-trivial topological lines, i.e. $\TVect_G(\Vect_{g_1},\Vect_{g_2})=\delta_{g_1,g_2}\Vect$.

Similarly as in the case of the fusion category $\Vect_G$, the $\TVect_G$ fusion rules can be twisted by a choice of non-trivial pentagonator that can be expressed in terms of a 4-cocycle $[\pi]\in H^4(G,\rU(1))$. As such, one obtains the 2-category denoted $\TVect_G^\pi$.
\subsubsection{Finite group 2-representation: $\TRep(G)$} The fusion 2-category $\TRep(G)$ admits a few equivalent definitions. For one, it can be understood as the \emph{Morita dual} of $\TVect_G$ with respect to the \emph{module 2-category} $\TVect$~\cite{Delcamp2021}, so that it can be realised as the fusion 2-category of \emph{$\TVect_G$-module 2-endofunctors} of $\TVect$, viz. $\TRep(G)=\TFun_{\TVect_G}(\TVect,\TVect)$. More useful for our discussion is that it can also be constructed as the category of \emph{$\Vect_G$-module categories and $\Vect_G$-module functors}. This is notated as $\TRep(G)\equiv\Mod(\Vect_G)$, in direct analogy with the fusion 1-categorical counterpart $\Rep(G)=\Mod(\mbb C[G])$. Elements of $\TRep(G)$ are therefore often said to encode \emph{$G$--2-representations}.

\bigskip\noindent Either way, simple objects are in one-to-one correspondence with the indecomposable $\Vect_G$-module categories and thus indexed by pairs $(H\subseteq G,\psi)$, $[\psi]\in H^2(H,\rU(1))$ (\cref{par:VecGMods}). We denote the associated defect by $\mbb U^{H_\psi}$. Its fusion rules are then specified by
\begin{equation}
	\mbb U^{H_\psi} \otimes \mbb U^{K_\varphi} = \bigoplus_{[r]\in H\backslash G / K} \mbb U^{(H\cap rKr^{-1})_{\varrho_r}}.
\end{equation}
Here, $\varrho_r$ is a 2-cocycle constructed from $\psi$ and $\varphi$ on the group $H\cap rKr^{-1}$, derived in ref.~\cite{Greenough2010}.  From these fusion rules we infer that the monoidal unit is given by $\mbb U^G$. The 1-morphisms between simple surfaces $\mbb U^{H_\psi}$, $\mbb U^{K_\varphi}$ are organised in the category of module functors
\begin{equation}
	\fr C(\mbb U^{H_\psi}, \mbb U^{K_\varphi}) = \Fun_{\Vect_G}(\mc R(H,\psi), \mc R(K,\varphi)),
\end{equation}
which concretely is given by
\begin{equation}\label{eq:2RepG1morph}
	\Fun_{\Vect_G}(\mc R(H,\psi), \mc R(K,\varphi)) = \bigoplus_{[r]\in H\backslash G / K} \Rep^{\varpi_r} (H\cap rKr^{-1}),
\end{equation}
where $\varpi_r$ is a 2-cocycle on $H\cap rKr^{-1}$ whose explicit expression in terms of $\psi$ and $\varphi$ can be found in ref.~\cite{Ostrik2001}. In particular, this illustrates the fact that self-interfaces on the identity surface are encoded in the symmetrically braided fusion category $(\Vect_G)^\star_\Vect=\Rep(G)$. These are said to encode the \emph{Wilson loop} operators, generalising those discussed above in the context of the toric code to general groups $G$. At the other extreme, one has $(\Vect_G)^\star_{\Vect_G}=\Vect_G$, which encodes the \emph{'t Hooft lines}.

\bigskip\noindent Despite their Morita equivalence, the structure of $\TVect_G$ and $\TRep(G)$ differs significantly. Clearly, $\TVect_G$ and $\TRep(G)$ are not monoidally equivalent, not even when $G$ is abelian. This highlights once more a remarkable distinction with the 1-categorical setting where $\Vect_G\simeq\Rep(G)$ as fusion categories, for abelian $G$.

Note that any two simple objects in $\TRep(G)$ belong to the same connected component, as we can conclude from inspecting \cref{eq:2RepG1morph}. In particular, all simple objects are connected to the transparent surface, which implies that all surfaces can be obtained via condensation of an appropriate set of Wilson loops encoded in $\Rep(G)$. More precisely, as argued in \cref{app:TwistedGauging}, the appropriate set of loops to condense is encoded in an \emph{algebra object} of $\Rep(G)$. In brief, an algebra object $A$ is a collection of lines $A=\bigoplus_\rho N_{\rho}^A \, \rho$ together with a \emph{multiplication map} $m:A\otimes A\mapsto A$, whose main property is that its composition is associative on the nose. We refer to \cref{app:gauging} for a more precise definition of algebra objects in unitary fusion categories. Provided an algebra object in $\Rep(G)$, the condensation defect thus schematically boils down to the following network of defects:
\begin{equation}\label{eq:CD_lines}
	\inputtikz{CD_lines}\,\,,
\end{equation}
that formally lives on the identity surface. We note that a representative algebra object associated with a surface $\mbb U^H$ corresponds, at the level of the objects, to the \emph{induced representation} ${\rm Ind}_H^G(\ub 1_H)$ in $G$ of the trivial representation $\ub 1_H$ of $H$~\cite{Davydov2010}. The associated multiplication morphism admits a simple expression in terms of the $\Rep(H)$ module associator over $\Rep(G)$~\cite{Fuchs2002,Bhardwaj2017,Vancraeynest2025b}. For the case of $G=\mbb Z_2$, the algebra object $\mbb C[\mbb Z_2]\simeq {\rm Ind}_{\{\mbb 1\}}^{\mbb Z_2}(\ub 1_{\{\mbb 1\}})$, viewed as the \emph{regular} $\mbb Z_2$-representation, results in the condensation defect \cref{eq:CD} upon expressing \cref{eq:CD_lines} in terms of the associated MPO symmetries.

Conversely, the surface $\mbb U^{H_\psi}$ can be constructed by condensing on $\mbb U^{\{\mbb 1\}}$ the algebra object $\mbb C^\psi[H]$ in the category $\Vect_G$ of 't Hooft lines, viewed as the $\psi$-twisted group algebra of $H$. As an object, $\mbb C^\psi[H]$ decomposes as $\bigoplus_{h\in H} h$, with the multiplication map $m$ implementing the $H$ group multiplication twisted by the 2-cocycle $\psi$.
\section{Tensor network representations of 2-categorical symmetries and dualities}\label{eq:2dDualities}
Having sketched the algebraic structure of fusion 2-categorical symmetries, with particular emphasis on $\TVect_G$ and $\TRep(G)$, we now turn in this concluding section to their lattice realisations in terms of PEPOs and intertwining MPOs, as well as the dualities associated with these symmetry structures~\cite{Haegeman2014a,Delcamp2021,Delcamp2023,Vancraeynest2025a}. As one might expect, the two-dimensional framework closely parallels the one-dimensional story told above. Owing to the relative novelty of this framework and technical intricacies of the \emph{Morita theory} of fusion 2-categories, this section remains largely qualitative~\cite{Decoppet2023b}.

For concreteness, let us focus on the case in which the input fusion 2-category is $\TVect_G$. As in one dimension, the input fusion 2-category determines an algebra of local symmetric operators from which a local symmetric Hamiltonian can be constructed. As demonstrated in refs.~\cite{Delcamp2021,Delcamp2023}, this algebra can be realised on the lattice via a choice of a \emph{module 2-category}. Such a choice determines in turn a \emph{Morita dual} fusion 2-category describing the symmetry structure of the resulting model. Consequently, the module 2-category naturally has the structure of a \emph{bimodule 2-category}.

One class of indecomposable module 2-categories is in one-to-one correspondence with pairs $(H\subseteq G,\omega)$ with $[\omega]\in H^3(H,\rU(1))$. As glossed over above, the Morita dual is defined as the fusion 2-category of \emph{module 2-endofunctors and module 2-natural transformations}. These 2-endofunctors and 2-natural transformations are equipped with coherence data, which in turn determine the components of the PEPO and MPO tensors supported on the corresponding surfaces. In many cases of interest, and in particular for $\TVect$ as $(\TVect_G,\TRep(G))$-bimodule 2-category, this construction simplifies substantially and admits a more explicit description, as illustrated in the aforementioned references.

Duality transformations amount to a change in the module 2-category, while keeping the input $\TVect_G$ fixed. Indeed, this change can be shown to leave invariant the algebra of local symmetric operators. The corresponding intertwining PEPOs implementing the duality are then naturally described in terms of \emph{$\TVect_G$-module 2-functors}. As in the one-dimensional setting, the construction simplifies considerably when one of the module 2-categories is chosen to be the regular one.

As an illustrative example, let us finally consider the Kramers-Wannier duality, which is realised by the change of module 2-category $\TVect_{\mbb Z_2}\mapsto\TVect$ over $\TVect_{\mbb Z_2}$. By specialising the general framework to this case and simplifying the duality operator, one recovers precisely the PEPO defined in \cref{eq:KWPEPO}.
